\documentclass[10pt,a4paper]{amsart}
\usepackage[margin=19mm]{geometry}
\usepackage{amsmath,amssymb,mathtools}
\usepackage[scr=rsfs,cal=euler]{mathalfa}
\usepackage{xfrac}
\usepackage[unicode,hidelinks,bookmarks=false]{hyperref}
\newcommand{\ie}{i.e.\ }
\newcommand{\cf}{cf.\ }
\newcommand{\resp}{respectively\ }
\newcommand{\Subsection}{\subsection}
\providecommand{\nobreakdash}{\nobreak\hbox{-}}
\setlength{\emergencystretch}{2em}
\multlinegap0pt
\usepackage{amssymb}
\usepackage{bm}
\usepackage{xcolor}
\usepackage{float}
\usepackage[shortlabels]{enumitem}
\newtheorem{lemma}{Lemma}
\newtheorem{proposition}{Proposition}
\title{Uniqueness and stability in bottom detection through surface measurements of water waves}
\author{Noureddine Lamsahel, Lionel Rosier}
\date{Source: arXiv:2601.17639v2 (CC BY 4.0)}
\begin{document}
\maketitle

\noindent\emph{Source and licence.} Uniqueness and stability in bottom detection through surface measurements of water waves. Version arXiv:2601.17639v2 (CC BY 4.0), distributed by arXiv under the Creative Commons Attribution 4.0 International licence, http://creativecommons.org/licenses/by/4.0/, as recorded in the arXiv licence metadata for this exact version and shown on the abstract page https://arxiv.org/abs/2601.17639v2. Full licence notice: \texttt{licenses/arxiv-2601.17639-CC-BY-4.0.txt}.\par\smallskip

\noindent\textit{Editorial scope.} Counted unit: the article's proposition lemmaaboutupperboundofinte (source0.tex lines 817--860). The article's complete original proof of it is delivered with it (source0.tex lines 861--863, one \texttt{proof} environment of 3 lines). No mathematics is written, repaired or re-derived: the statement and the proof are the article's own lines, in the article's own notation, and no retained line is substituted or re-flowed.  Retained uncounted as the source-local context the proof needs: lines 278--280; lines 283--294; lines 300--311; lines 403--405; lines 469--471; lines 737--754; lines 782--813; lines 792--794; lines 1187--1187.  Nothing was omitted from any retained span, so the package carries no omission record.  No retained line contains a citation, so the wrapper carries no bibliography and no bibliography entry.  No retained line contains a figure environment, an image-inclusion command or drawing code, so no image is delivered and the package declares no asset dependency.  Editorial formatting only, in addition: an AMS \texttt{amsart} preamble that restates the article's own declarations and macros used by the retained lines, namely the environments \texttt{lemma}, \texttt{proposition} on separate counters, together with the standard packages those lines need.  The retained environments print in this document as Lemma 1, Proposition 1 in order of appearance, whereas the article prints the same items with the numbering of its own full document.  The complete native source of the article as distributed in the arXiv e-print for this exact version is bundled unchanged as \texttt{sources/193335/source0.tex}.
\begin{equation}\label{conditionofliquid}
	\zeta_0(X)-b_0(X)\geq H_0\;\;\text{and }\;\; \zeta(X)-b(X)\geq H_0,\quad \forall X\in \mathbb{R} ^d. 
\end{equation}

\begin{equation}\label{pd3}
	\left\{
	\begin{array}{rrrrr}
		\vspace{0.1cm}
		\Delta\phi &=&0,&\quad \text{on}\; \; \Omega_{b}^{\zeta},\\
		\vspace{0.1cm}
		\phi &=&\psi,&\quad \text{in}\;\;  \Gamma^{\zeta},\\
		\partial_n\phi &=&0,&\quad \text{in}\;\;\Gamma^{b},\\
		\phi &=&\theta ,&\quad \text{in}\;\;  \Gamma^{\zeta}_b,\\
	\end{array}
	\right.
\end{equation}

\begin{equation}\label{pd4}
	\left\{
	\begin{array}{rrrrr}
		\vspace{0.1cm}
		\Delta\phi_0 &=&0,&\quad \text{on}\; \; \Omega_{b_0}^{\zeta_0},\\
		\vspace{0.1cm}
		\phi_0 &=&\psi_0,&\quad \text{in}\;\;  \Gamma^{\zeta_0},\\
		\partial_n\phi_0 &=&0,&\quad \text{in}\;\;\Gamma^{b_0},\\
		\phi _0 &=&\theta _0 ,&\quad \text{in}\;\;  \Gamma^{\zeta_0}_{b_0},\\
	\end{array}
	\right.
\end{equation}

\begin{equation}\label{defs1ands2}
	S_1=(\zeta-\zeta_0)^{-1}(\mathbb{R}^*_{+}),\;\;S_2=(\zeta-\zeta_0)^{-1}(\mathbb{R}_{-}),
\end{equation}

	\begin{equation}\label{conditiononsurfaces}
		|\zeta(X)-\zeta_0(X)|\leq \frac{H_0}{2},\quad \forall X\in \mathcal O .  
	\end{equation}

\begin{lemma}\label{l1stab}
	Let $b$, $b_0$, $\zeta$ and $\zeta_0$ be  functions in $C^1\left(\overline{\mathcal O }\right)$ such that (\ref{conditionofliquid}) and (\ref{conditiononsurfaces}) hold. Let 
	$\phi\in H ^2(\Omega_b^{\zeta})$ be a solution of system (\ref{pd3}). Then 
	\begin{equation}\label{upperforsurfacephi}
		\begin{split}
			\left\Vert  \psi-\phi_{\zeta_0}   \right\Vert_{L^2(S_1)}&\leq\left\Vert \zeta-\zeta_0  \right\Vert^{\frac{1}{2}}_{L^{\infty}(\mathcal O)}\left\Vert  \phi  \right\Vert_{H ^2(\Omega_b^{\zeta})},\\
			\left\Vert  (\nabla_X\phi )_{\zeta}-(\nabla_X\phi )_{\zeta_0}   \right\Vert_{L^2(S_1)}&\leq\left\Vert \zeta-\zeta_0  \right\Vert^{\frac{1}{2}}_{L^{\infty}(\mathcal O)}\left\Vert  \phi  \right\Vert_{H ^2(\Omega_b^{\zeta})},\\
			\left\Vert  (\partial_y\phi )_{\zeta}-(\partial_y\phi )_{\zeta_0}   \right\Vert_{L^2(S_1)}&\leq\left\Vert \zeta-\zeta_0  \right\Vert^{\frac{1}{2}}_{L^{\infty}(\mathcal O)}\left\Vert  \phi  \right\Vert_{H ^2(\Omega_b^{\zeta})},
		\end{split}
	\end{equation} 
	where $\phi_{\zeta_0}$, $(\nabla_X \phi)_{\zeta _0}$ and $(\partial _y\phi ) _{\zeta _0}$   are defined by 
	\begin{equation}\label{defforsimplication}
		\phi_{\zeta_0}(X)=\phi(X,\zeta_0(X)),\;\; (\nabla_X\phi )_{\zeta _0}(X)= (\nabla_X\phi ) (X,\zeta _0(X)),\;\; (\partial _y\phi )_{\zeta _0}(X)= 
		(\partial _y\phi ) (X,\zeta_0(X)),\;\forall X\in S_1,
	\end{equation}
	$(\nabla_X\phi)_{\zeta }$ and $(\partial _y\phi ) _{\zeta }$ being defined similarly, 
	and the set $S_1$ is given by (\ref{defs1ands2}).
\end{lemma}

\begin{lemma}\label{uppergradfortheupper}
	Let $b$, $b_0$, $\zeta$ and $\zeta_0$ be functions in $C^1\left(\overline{\mathcal O}\right)$ such that (\ref{conditionofliquid}) and (\ref{conditiononsurfaces}) hold, and  let $\phi\in H^2(\Omega^{\zeta}_{b})$, $\phi_0\in H^2(\Omega^{\zeta_0}_{b_0})$ be solutions of systems (\ref{pd3}) and (\ref{pd4}), respectively. Then
	\begin{equation}\label{relationbetweengradandourmesurement}
		\begin{split}
			\left\Vert (\nabla_X\phi )_{\zeta}-(\nabla_X\phi_0)_{\zeta_0}  \right\Vert_{L^2( \mathcal O )}&\leq \left[G_1\left(\zeta-\zeta_0,\psi-\psi_0,\partial_n\phi_{|_{\Gamma^{\zeta}}}-\partial_n\phi_{0|_{\Gamma^{\zeta_0}}}\right)\right]^{1/2},\\
			\left\Vert (\partial_y\phi ) _{\zeta}- (\partial_y\phi_0)_{ \zeta_0}  \right\Vert_{L^2(\mathcal O )}&\leq \left[G_1\left(\zeta-\zeta_0,\psi-\psi_0,\partial_n\phi_{|_{\Gamma^{\zeta}}}-\partial_n\phi_{0|_{\Gamma^{\zeta_0}}}\right)\right]^{1/2},
		\end{split}
	\end{equation}
	where $(\nabla_X\phi )_{\zeta},  (\partial _y\phi )_{\zeta}, (\nabla_X\phi _0)_{\zeta _0}, 
	(\partial _y\phi_0 )_{\zeta_0}$ are defined according to (\ref{defforsimplication}). The term $G_1$  is given by
	  \begin{equation}\label{firstmajG1}
	  	G_1=\sup(\widehat{G}_1,\widetilde{G}_1),
	  \end{equation} 
	  where, 
	\begin{equation}\label{firstmajG11}
		\begin{split}
			\widehat{G}_1(\zeta-\zeta_0, \psi-\psi_0 ,\partial_n\phi_{|_{\Gamma^{\zeta  }     }}-\partial_n\phi_{0|_{\Gamma^{\zeta_0    }   }} )&= 3\left\Vert \nabla_X\psi-\nabla_X\psi_0  \right\Vert_{L^{2}(\mathcal O )}^2+3\left\Vert  \partial_n\phi_{|_{\Gamma^{\zeta   }    }}-\partial_n\phi_{0|_{\Gamma^{\zeta_0 }      }}\right\Vert_{L^{2}(\mathcal O )}^2\\
		&\hspace{0.3cm}+ 3\left\Vert z_3 \right\Vert_{L^{2}(\mathcal O )}^2 \left\Vert \nabla_X\zeta_0-\nabla_X\zeta   \right\Vert_{L^{\infty}(\mathcal O )}^2,\\
		\end{split}
	\end{equation} 
	and 
	\begin{equation}\label{firstmajG12}
		\begin{split}
			\widetilde{G}_1(\zeta-\zeta_0, \psi-\psi_0 ,\partial_n\phi_{|_{\Gamma^{\zeta  }     }}-\partial_n\phi_{0|_{\Gamma^{\zeta_0    }   }} )&=3\left\Vert \nabla_X\psi-\nabla_X\psi_0  \right\Vert_{L^{2}(\mathcal O )}^2+3\left\Vert \nabla_X\zeta_0 \right\Vert_{L^{\infty}(\mathcal O )}^2\widehat{G}_1\\
			&\hspace{0.3cm}+3 \left\Vert (\partial_y\phi)_{\zeta} \right\Vert_{L^{2}(\mathcal O )}^2 \left\Vert \nabla_X\zeta_0-\nabla_X\zeta   \right\Vert_{L^{\infty}(\mathcal O )}^2.\\
		\end{split}
	\end{equation}
	In  \eqref{firstmajG11}, the coefficient $z_3$ is given  by 
	\begin{equation}\label{z1z2z3}
			z_3=\left| \partial_n\phi_{0|_{\Gamma^{\zeta_0       }}} \right|+ \left|\nabla_X\psi_0  \right|+(1+\left| \nabla_X\zeta_0  \right|)\left|(\partial_y\phi_0)_{\zeta_0}  \right|.
	\end{equation}
\end{lemma}

	  \begin{equation}
	  	G_1=\sup(\widehat{G}_1,\widetilde{G}_1),
	  \end{equation} 

\section{Proof of Proposition \ref{lemmaaboutupperboundofinte} }\label{A.2}

\begin{proposition}\label{lemmaaboutupperboundofinte}
	Let $b$, $b_0$, $\zeta$ and $\zeta_0$ be functions in  $C^1\left(\overline{\mathcal O}\right)$ such that (\ref{conditionofliquid}) and (\ref{conditiononsurfaces}) hold, and  let  $\phi\in H^2(\Omega^{\zeta}_{b})$, $\phi_0\in H^2(\Omega^{\zeta_0}_{b_0})$ be solutions of systems (\ref{pd3}) and (\ref{pd4}), respectively. Then
	
	\begin{equation}\label{upperboundsforintegral}
		\begin{array}{ll}
			\vspace{0.2cm}
			\displaystyle\int_{\Gamma^{\underline{\zeta}}} \partial_n\phi(\phi-\phi_0)&\hspace{-0.7cm}\leq G_2\left(\zeta-\zeta_0,\psi-\psi_0\right),\\
			\vspace{0.2cm}
			\displaystyle\int_{\Gamma^{\underline{\zeta}}} \phi_0(\partial_n\phi_0-\partial_n\phi)&\hspace{-0.2cm}\leq G_3\left(\zeta-\zeta_0,\partial_n\phi_{|_{\Gamma^{\zeta}}}-\partial_n\phi_{0|_{\Gamma^{\zeta_0}}}\right),\\
			\vspace{0.2cm}
			||\phi-\phi_0||_{L^2(\Gamma^{\underline{\zeta}})}&\hspace{-1cm}\leq \left[G_4\left(\zeta-\zeta_0,\psi-\psi_0\right)\right]^{1/2},\\
			||\nabla (\phi-\phi_0)||_{L^2(\Gamma^{\underline{\zeta}})}&\hspace{-0.5cm}\leq \left[G_5\left(\zeta-\zeta_0,\psi-\psi_0,\partial_n\phi_{|_{\Gamma^{\zeta}}}-\partial_n\phi_{0|_{\Gamma^{\zeta_0}}}\right)\right]^{1/2},
		\end{array}    
	\end{equation}
	where  $G_2$, $G_3$, $G_4$ and $G_5$ are given by
	\begin{equation}\label{G2to4exprsl}
		\begin{array}{ll}
			G_2(\zeta-\zeta_0,\psi-\psi_0)&\hspace{-3cm} :=\left[\left\Vert   z_4\right\Vert_{L^2(S_1)}\left\Vert\phi\right\Vert_{H^2(\Omega_{b}^{\zeta})} +\left\Vert   z_5\right\Vert_{L^2(\mathcal O)}\left\Vert\phi_0\right\Vert_{H^2(\Omega_{b_0}^{\zeta_0})} \right] \left\Vert \zeta-\zeta_0  \right\Vert_{L^{\infty}(\mathcal O )}^{\frac{1}{2}}\\\vspace{0.3cm}
			&\hspace{0.5cm}+\left[\left\Vert   z_4\right\Vert_{L^2(S_1)} +\left\Vert   z_5\right\Vert_{L^2(\mathcal O)} \right] \left\Vert   \psi-\psi_0 \right\Vert_{L^2(\mathcal O )},\\
			G_3\left(\zeta-\zeta_0,\partial_n\phi_{|_{\Gamma^{\zeta}}}-\partial_n\phi_{0|_{\Gamma^{\zeta_0}}}\right)&\hspace{-1.4cm}:=z_7\left[\left\Vert\phi_{0\zeta_0}\right\Vert_{L^2(\mathcal O)}+\left\Vert\phi_{0\zeta}\right\Vert_{L^2(S_2)}\right] \left\Vert\partial_n\phi_{|_{\Gamma^{\zeta}}}-\partial_n\phi_{0|_{\Gamma^{\zeta_0}}}\right\Vert_{L^2(\mathcal O )}\\
			&\hspace{-0.35cm}+z_7\Bigg\{  \left(  \left\Vert \phi_{0\zeta_0} \right\Vert_{L^2(\mathcal O)}+ \left\Vert  \phi_{0\zeta_0}\nabla_X\zeta_0\right\Vert_{L^2(\mathcal O)} \right)\left\Vert \phi  \right\Vert_{H^2(\Omega_b^{\zeta})}  \\
			&\hspace{-0.3cm}+ \left(  \left\Vert \phi_{0\zeta} \right\Vert_{L^2(S_2)}+ \left\Vert  \phi_{0\zeta} \nabla_X\zeta_0\right\Vert_{L^2(S_2)} \right)\left\Vert \phi_0  \right\Vert_{H^2(\Omega_{b_0}^{\zeta_0})}      \Bigg\}\left\Vert   \zeta-\zeta_0\right\Vert_{L^{\infty}(\mathcal O )}^{\frac{1}{2}}\\\vspace{0.3cm}
			&\hspace{-0.3cm}+z_7z_6\left\Vert\nabla_X\zeta-\nabla_X\zeta_0  \right\Vert_{L^{\infty}(\mathcal O )},\\
		 G_4(\zeta-\zeta_0,\psi-\psi_0)&\hspace{-3cm} :=z_7\bigg[4 \left\Vert \psi-\psi_0 \right\Vert_{L^{2}(\mathcal O )}^2
		+  2 (\left\Vert\phi\right\Vert_{H^2(\Omega_b^{\zeta})}^2+\left\Vert\phi_0\right\Vert_{H^2(\Omega_{b_0}^{\zeta_0})}^2)\left\Vert\zeta-\zeta_0\right\Vert_{L^{\infty}(\mathcal O )} \bigg],  \\   \vspace{0.3cm}
		G_5\left(\zeta-\zeta_0,\psi-\psi_0,\partial_n\phi_{|_{\Gamma^{\zeta}}}-\partial_n\phi_{0|_{\Gamma^{\zeta_0}}}\right)&\hspace{-0.2cm} 
		:=4z_7\left[  \left\Vert  \phi \right\Vert_{H^2(\Omega_{b}^{\zeta})}^2+   \left\Vert  \phi_0 \right\Vert_{H^2(\Omega_{b_0}^{\zeta_0})}^2 \right]\left\Vert   \zeta-\zeta_0\right\Vert_{L^{\infty}(\mathcal O )}  +8z_7 G_1.
		\end{array}
	\end{equation}
	
	Here, the term  $G_1$ is defined in (\ref{firstmajG1}) and  coefficients $z_4$, $z_5$, $z_6$, and $z_7$ are given by 
	\begin{equation}\label{z4toz9}
		\left\{
		\begin{array}{ll}
			z_4 &:= \sqrt{1+|\nabla_X\zeta_0|^2}\;\partial_n\phi_{|_{\Gamma^{\zeta_0}}}, \\
			z_5 &:= \sqrt{1+|\nabla_X\zeta|^2}\;\partial_n\phi_{|_{\Gamma^{\zeta}}},  \\
			z_6&:=  \left\Vert \phi_{0\zeta} \right\Vert_{L^2(S_2)}\left[  \left\Vert  (\nabla_X\phi_0)_{\zeta}\right\Vert_{L^2(S_2)}+ \left\Vert \partial_n\phi_{0|_{\Gamma^{\zeta_0}}} \right\Vert_{L^2(\mathcal O)} \right],\\
			&\hspace{0.5cm}+\left\Vert \phi_{0\zeta_0} \right\Vert_{L^2(\mathcal O)}\left[  \left\Vert  (\nabla_X\phi)_{\zeta}\right\Vert_{L^2(\mathcal O)}+ \left\Vert \partial_n\phi_{|_{\Gamma^{\zeta_0}}} \right\Vert_{L^2(S_1)} \right],\\
			z_7&:=   \max\left\{\left\Vert \sqrt{1+|\nabla_X\zeta|^2} \right\Vert_{L^{\infty}(\mathcal O)},\left\Vert \sqrt{1+|\nabla_X\zeta_0|^2} \right\Vert_{L^{\infty}(\mathcal O)}  \right\}.
		\end{array}
		\right.
	\end{equation}
	
\end{proposition}

\begin{proof} Let us sketch the proof. 
	We transform the surface integrals on the left-hand side of (\ref{upperboundsforintegral}) into Lebesgue integrals over $\mathcal O$. Then, employing arguments similar to those used to prove  Lemma \ref{uppergradfortheupper}, we leverage the inequalities provided by Lemma \ref{l1stab} and Lemma \ref{uppergradfortheupper} to derive the desired upper bounds. For ease of reading, the proof is shifted to Appendix~\ref{A.2}.
\end{proof}

\end{document}
